\section{Moving between position, velocity, and acceleration}

\subsection{The derivative direction}

So far we have moved in one direction:
\[
\mathbf{r}(t)
\quad \longrightarrow \quad
\mathbf{v}(t)
\quad \longrightarrow \quad
\mathbf{a}(t).
\]
This direction uses derivatives:
\[
\mathbf{v}(t)=\frac{d\mathbf{r}}{dt},
\qquad
\mathbf{a}(t)=\frac{d\mathbf{v}}{dt}=\frac{d^2\mathbf{r}}{dt^2}.
\]

If we know position as a function of time, we can compute velocity and acceleration by differentiating. No additional information is needed for this direction. The function \(\mathbf{r}(t)\) already contains all information about the motion.

\begin{examplebox}
If
\[
\mathbf{r}(t)=(t,t^2),
\]
then
\[
\mathbf{v}(t)=(1,2t),
\qquad
\mathbf{a}(t)=(0,2).
\]
Everything follows by differentiation.
\end{examplebox}

\subsection{The integral direction}

We can also move in the opposite direction. If we know acceleration, we may want to recover velocity. If we know velocity, we may want to recover position:
\[
\mathbf{a}(t)
\quad \longrightarrow \quad
\mathbf{v}(t)
\quad \longrightarrow \quad
\mathbf{r}(t).
\]
This direction uses integration. Integration introduces constants or, equivalently, requires initial conditions. These are not mathematical noise; they represent missing information about the motion.

\begin{theorembox}
If \(\mathbf{v}(t)=d\mathbf{r}/dt\) on an interval containing \([t_0,t]\), then
\[
\mathbf{r}(t)=\mathbf{r}(t_0)+\int_{t_0}^{t}\mathbf{v}(\tau)\,d\tau.
\]
If \(\mathbf{a}(t)=d\mathbf{v}/dt\) on an interval containing \([t_0,t]\), then
\[
\mathbf{v}(t)=\mathbf{v}(t_0)+\int_{t_0}^{t}\mathbf{a}(\tau)\,d\tau.
\]
\end{theorembox}

The symbol \(\tau\) is a dummy variable of integration. The upper limit \(t\) is the time at which we want to find the position or velocity. The lower limit \(t_0\) is the initial time.

\subsection{Why initial conditions are necessary}

Knowing how something changes is not the same as knowing its value. Suppose we know only that
\[
v(t)=3.
\]
Then
\[
x(t)=\int 3\,dt=3t+C.
\]
The constant \(C\) is not determined by the velocity. It tells us the initial position. If \(x(0)=2\), then \(C=2\), so
\[
x(t)=3t+2.
\]

\begin{examplebox}
Suppose acceleration is constant:
\[
a(t)=A.
\]
Then
\[
v(t)=At+C_1.
\]
If \(v(0)=v_0\), then \(C_1=v_0\), so
\[
v(t)=v_0+At.
\]
Integrating again gives
\[
x(t)=v_0t+\frac{1}{2}At^2+C_2.
\]
If \(x(0)=x_0\), then \(C_2=x_0\), and therefore
\[
x(t)=x_0+v_0t+\frac{1}{2}At^2.
\]
\end{examplebox}

This formula is often presented as a standard kinematic formula. Here it appears as the result of integrating constant acceleration twice and using initial conditions.

\subsection{Vector form}

The same idea works in vector form. If acceleration is known, then
\[
\mathbf{v}(t)
=
\mathbf{v}(t_0)+\int_{t_0}^{t}\mathbf{a}(\tau)\,d\tau.
\]
If velocity is known, then
\[
\mathbf{r}(t)
=
\mathbf{r}(t_0)+\int_{t_0}^{t}\mathbf{v}(\tau)\,d\tau.
\]

If acceleration is constant and equal to \(\mathbf{A}\), and if \(t_0=0\), then
\[
\mathbf{v}(t)=\mathbf{v}_0+\mathbf{A}t,
\]
and
\[
\mathbf{r}(t)=\mathbf{r}_0+\mathbf{v}_0t+\frac{1}{2}\mathbf{A}t^2.
\]

\subsection{The central structure of kinematics}

The central structure of kinematics is
\[
\boxed{\mathbf{v}(t)=\frac{d\mathbf{r}}{dt}}
\]
and
\[
\boxed{\mathbf{a}(t)=\frac{d\mathbf{v}}{dt}=\frac{d^2\mathbf{r}}{dt^2}}.
\]
The inverse relations are
\[
\boxed{
\mathbf{v}(t)
=
\mathbf{v}(t_0)+\int_{t_0}^{t}\mathbf{a}(\tau)\,d\tau
}
\]
and
\[
\boxed{
\mathbf{r}(t)
=
\mathbf{r}(t_0)+\int_{t_0}^{t}\mathbf{v}(\tau)\,d\tau
}.
\]

These formulas should not be treated as isolated recipes. They express the logic of kinematics: position describes where the body is, velocity describes how position changes, acceleration describes how velocity changes, and initial conditions supply the information needed when reconstructing motion by integration.

Dynamics will ask why acceleration appears. Before that question can be answered, we need the language developed in this chapter: position, velocity, acceleration, derivatives, integrals, and initial conditions.
